\documentclass[10pt,a4paper]{amsart}
\usepackage[margin=19mm]{geometry}
\usepackage{amsmath,amssymb,mathtools}
\usepackage[scr=rsfs,cal=euler]{mathalfa}
\usepackage{xfrac}
\usepackage[unicode,hidelinks,bookmarks=false]{hyperref}
\newcommand{\ie}{i.e.\ }
\newcommand{\cf}{cf.\ }
\newcommand{\resp}{respectively\ }
\newcommand{\Subsection}{\subsection}
\providecommand{\nobreakdash}{\nobreak\hbox{-}}
\setlength{\emergencystretch}{2em}
\multlinegap0pt
\usepackage{bm}
\usepackage{bbm}
\usepackage{dsfont}
\usepackage{xspace}
\usepackage{cancel}
\usepackage{mathtools}
\usepackage{amsthm}
\makeatletter
\@ifundefined{theorem}{\newtheorem{theorem}{Theorem}}{}
\@ifundefined{lemma}{\newtheorem{lemma}[theorem]{Lemma}}{}
\@ifundefined{proposition}{\newtheorem{proposition}[theorem]{Proposition}}{}
\makeatother
\def \del {\partial} 
\title[H\"older regularity of the pressure for weak solutions of th]{H\"older regularity of the pressure for weak solutions of the 3D Euler equations in bounded domains}
\author{Claude Bardos, Daniel W. Boutros, Edriss S. Titi}
\date{Source: arXiv:2304.01952v2 (CC BY 4.0)}
\begin{document}
\maketitle

\noindent\emph{Source and licence.} H\"older regularity of the pressure for weak solutions of the 3D Euler equations in bounded domains. Version arXiv:2304.01952v2 (CC BY 4.0), distributed under the Creative Commons Attribution 4.0 International licence, as recorded in the arXiv licence metadata for this exact version and shown on the abstract page https://arxiv.org/abs/2304.01952v2. Full rights notice: licenses/arxiv-2304.01952-CC-BY-4.0.txt.\par\smallskip

\noindent\textit{Editorial scope.} Counted unit: the Proposition whose source label is \texttt{th1} in the article (source0.tex lines 1391--1427, source label \texttt{th1}), delivered together with the article's own complete original proof of it (source0.tex lines 1428--1528, one \texttt{proof} environment of 101 lines). No mathematics is written, repaired or re-derived: statement and proof are the article's own text line for line. Required context retained uncounted: veryweakboundcondeq (lines 203--205); usualboundarycond (lines 207--209); equivalencelemma (lines 215--221); basica1 (lines 1382--1388). Every definition, display and label of those lines is retained unchanged. Omitted from this package and not redistributed: 1--202, 206--206, 210--214, 222--1381, 1389--1390, 1529--1986. Nothing inside a retained excerpt is omitted or altered: every retained body is delivered exactly as the article's lines are written, so no omission receipt of any kind accompanies this package. Citations: the retained excerpts carry no citation command at all, so no bibliography entry of the article is transcribed and no reference is redistributed. Reference adaptation: none was needed and none was made; every reference inside the delivered unit resolves inside it, so no unresolved reference or citation is produced. Printed slips kept literal: no printed word, symbol, hypothesis or number of the counted statement or of its proof was altered. Layout and class: the article's own class is \texttt{article} with packages including graphicx, subfig, amsmath, amsthm, amssymb, esint, mathrsfs, appendix, hyperref, geometry, color, array, url, float; the wrapper is the lane's pinned \texttt{amsart} editorial wrapper at 10pt on A4, which restates the theorem-like environments the retained text needs and adds the article's own macro definitions that the retained text needs (\texttt{\textbackslash del}), with extra wrapper packages (bm, bbm, dsfont, xspace, cancel, mathtools). The delivered numbering is the unit's own. Assets: no retained span contains a figure environment, a graphics-inclusion command, a PDF-page inclusion, a table or a TikZ picture; no image file is copied into this package, the package declares no asset dependency and no image file is redistributed with it. Licence: version arXiv:2304.01952v2 of this article is distributed under the Creative Commons Attribution 4.0 International licence, as recorded in the arXiv licence metadata for this exact version and shown on the abstract page https://arxiv.org/abs/2304.01952v2. A full rights notice is delivered at licenses/arxiv-2304.01952-CC-BY-4.0.txt. Characters: every retained excerpt is pure ASCII, so the delivered engine, XeLaTeX with its default Latin Modern text font, reports no missing character.
\begin{equation} \label{veryweakboundcondeq}
\partial_n \big( p + (u \cdot n)^2 \big) = \big( u \otimes u \big) : \nabla n, \quad \text{on } \partial \Omega,
\end{equation}

\begin{equation} \label{usualboundarycond}
\partial_n p = \big( u \otimes u \big) : \nabla n, \quad \text{on } \partial \Omega.
\end{equation}

\begin{lemma} \label{equivalencelemma}
Let $u \in C^{0,\alpha} (\Omega)$ with $\alpha > \frac{1}{2}$ be a velocity field such that $(u \cdot n) \lvert_{\partial \Omega} = 0$, then
\begin{equation}
\partial_n (u \cdot n)^2 = 0,
\end{equation}
and consequently boundary conditions \eqref{veryweakboundcondeq} and \eqref{usualboundarycond} are equivalent.
\end{lemma}

\begin{equation}\label{basica1}
\begin{aligned}
&v(x,y) \coloneqq (u_2(x,y))^2= \sum_{k_1, k_2=0}^\infty \bigg( 2^{-\alpha (k_1 + k_2)} \cos(2^{k_1} \pi x)  \cos(2^{k_2}\pi  x)  \sin(2^{k_1}y)  \sin(2^{k_2} y)\bigg), \\
&U(y;\theta) =  \sum_{k_1, k_2=0}^\infty \int_{\mathbb{T}} v (x,y) \theta (x)dx \\
&= \sum_{k_1, k_2=0}^\infty \bigg( \int_{\mathbb{T}} 2^{-\alpha (k_1 + k_2)} \cos(2^{k_1}  \pi x)  \cos(2^{k_2}\pi  x)\theta (x) dx\bigg)  \sin(2^{k_1}y)  \sin(2^{k_2} y) .
\end{aligned}
\end{equation}

\begin{proposition}\label{th1}
$$\,$$
\begin{enumerate}
\item Suppose $\theta \in \mathcal D({\mathbb{T}})$ satisfies
\begin{equation}
\int_{\mathbb{T}} \theta(x) dx =0,
\end{equation}
then the limit
$$
\lim_{y\rightarrow 0 +} \frac1{y}  \int_{\mathbb{T}} v (x,y) \theta (x)dx
$$
is well-defined and is equal to $0\,.$

\item Otherwise if $0 < \alpha < \frac{1}{2}$ and
\begin{equation}
\int_{\mathbb{T}} \theta(x) dx \not=0,
\end{equation}
then it holds that
$$
\liminf _{y \rightarrow 0 +} \frac1{y } \bigg\lvert \int_{\mathbb{T}} v (x,y) \theta (x)dx \bigg\rvert =\infty.
$$
As a consequence the function
\begin{equation}
U(y; \theta) = \int_{\mathbb{T}} v (x,y) \theta (x)dx
\end{equation}
does not have a well-defined derivative at the point $y=0\,.$ 
\item If $\alpha = \frac{1}{2}$ and 
\begin{equation}
\int_{\mathbb{T}} \theta(x) dx \not=0,
\end{equation}
then we have that
\begin{equation}
\liminf _{y \rightarrow 0 +} \frac1{y } \bigg\lvert \int_{\mathbb{T}} v (x,y) \theta (x)dx \bigg\rvert \geq 2.
\end{equation}
Consequently, if the derivative of $U(y,\theta)$ exists it is not equal to zero.
\end{enumerate}
\end{proposition}

\begin{proof}
For the proof one first eliminates the nonresonant terms (i.e., the terms involving $k_1\not=k_2$ ) and then a comparison argument is used. As such the subscripts $R$ and $NR$ are used to denote the resonant and nonresonant parts of $U$, respectively.

First one has the following.
\begin{lemma}
The function
\begin{equation}
U_{NR} (y,\theta) \coloneqq \sum_{k_1,k_2 = 0,k_1\not=k_2}^\infty \bigg(2^{-\alpha (k_1 + k_2)} \bigg[ \int_{\mathbb{T}} \cos(2^{k_1}  \pi x )  \cos(2^{k_2}  \pi x)\theta (x) dx \bigg] \sin(2^{k_1}y)  \sin(2^{k_2} y)  \bigg)
\end{equation}
belongs to $C^1([0,1])$, moreover we have that $\del_y U_{NR} (y;\theta) \lvert_{y=0} =0$.
\end{lemma}
\begin{proof}
One first recalls the following trigonometric identity
\begin{equation*}
\cos(2^{k_1}  \pi x)  \cos(2^{k_2} \pi x)=\frac12 \bigg[\cos((2^{k_1}+ 2^{k_2})  \pi x) +\cos((2^{k_1}- 2^{k_2}) \pi x)\bigg].
\end{equation*}
Then for $k_1 \neq k_2$ and $m > 0$ it holds that
\begin{equation} \label{nonres}
\begin{aligned}
&\int_{\mathbb{T}} \cos(2^{k_1}  \pi x)  \cos(2^{k_2}  \pi x)\theta (x) dx=\frac{(-1)^{m}}{2( \pi( 2^{k_1}+ 2^{k_2}))^{2m} } \int_{\mathbb{T}}  \cos((2^{k_1}+ 2^{k_2}) \pi x) \frac{d^{2m}}{dx^{2m}}\theta(x) dx\\
&+\frac{(-1)^{m}}{2(\pi(2^{k_1}- 2^{k_2}))^{2m} } \int_{\mathbb{T}} \cos((2^{k_1}- 2^{k_2}) \pi x) \frac{d^{2m}}{dx^{2m}}\theta(x) dx,
\end{aligned}
\end{equation}
%\end{document}
moreover we have
\begin{equation}\label{nonres2}
\frac{d}{dy} ( \sin(2^{k_1} \pi y)  \sin(2^{k_2} \pi y) )= 2^{k_1} \pi \cos(2^{k_1} \pi y)  \sin(2^{k_2} \pi y) + 2^{k_2} \pi \sin(2^{k_1} \pi y)  \cos(2^{k_2} \pi y).
\end{equation}
Combining equations (\ref{nonres}) and (\ref{nonres2} ) one obtains that
\begin{equation}
\begin{aligned}
&\bigg\lvert\frac{d}{dy} \bigg[ \bigg( \int_{\mathbb{T}} \cos(2^{k_1}\pi x)  \cos(2^{k_2}\pi x)\theta (x) dx \bigg)  \sin(2^{k_1}y)  \sin(2^{k_2} y) \bigg] \bigg\rvert\\
&\le C \pi \bigg( \frac 1{2(\pi (2^{k_1}+ 2^{k_2}))^{2m} } +
   \frac 1{2((\pi(2^{k_1}- 2^{k_2}))^{2m} } \bigg) (2^{k_1} + 2^{k_2} )\int_{\mathbb{T}} \bigg\lvert\frac{d^{2m}}{dx^{2m}} \theta(x)\bigg\rvert dx 
\end{aligned}
 \end{equation}
which for $m \geq 1$ is a sequence constituting the terms of an absolutely converging series.

As a consequence  the function
\begin{equation*}
U_{NR}(y;\theta)= \sum_{k_1,k_2=0 , k_1\not=k_2}^\infty \bigg(2^{-\alpha (k_1 + k_2)} \bigg[ \int_{\mathbb{T}} \cos(2^{k_1}\pi x)  \cos(2^{k_2}\pi  x)\theta (x) dx \bigg] \sin(2^{k_1} \pi y)  \sin(2^{k_2} \pi y) \bigg)  
\end{equation*}
belongs to the space $C^{1}([0,1])$ and satisfies the relation:
\begin{equation}
\del_y U_{NR} (y;\theta) \lvert_{y=0} =0.
\end{equation}

\end{proof}
Therefore one only has to consider the resonant part of the Weierstrass series (i.e. the case $k_1 = k_2$ in equation \eqref{basica1}), which is
\begin{equation}
U_R(y;\theta) =  \sum_{k=0 }^\infty 2^{-2\alpha k} \bigg( \int_{\mathbb{T}} (\cos(2^{k } \pi  x)  )^2  \theta (x) dx \bigg)  (\sin(2^{k } \pi y))^2.
\end{equation}
By using the identity
\begin{equation}
 (\cos(2^{k } \pi x)  )^2=\frac12 \bigg(1+ \cos(2^{k+1 } \pi  x)  \bigg),
\end{equation}
one has that
\begin{align*}
U_R(y;\theta) &= \frac12 \bigg( \int_{\mathbb{T}} \theta (x) dx \bigg) \sum_{k =0}^\infty \bigg( 2^{-2\alpha k} (\sin(2^{k }\pi y) )^2 \bigg) \\
&+ \frac{1}{2} \sum_{k =0}^\infty \bigg[ 2^{-2\alpha k} \bigg( \int_{\mathbb{T}} \cos(2^{k +1}  \pi x)    \theta (x) dx \bigg) (\sin(2^{k } \pi y))^2 \bigg].
\end{align*}
Similarly as before, the function
\begin{equation}
U_{RNR}(y;\theta)= \frac{1}{2} \sum_{k=0}^\infty 2^{-2\alpha k} \bigg( \int_{\mathbb{T}} \cos(2^{k +1}\pi  x) \theta (x) dx \bigg)  (\sin(2^{k } \pi y))^2
\end{equation}
is a series converging in $C^{1}([0,1]) $ with derivative equal to $0$ for $y=0$.
Hence the completion of the proof now relies only on the analysis of the behaviour of the term
\begin{equation}
\frac{1}{y} U_{RR}(y;\theta)= \frac{1}{2y} \int_{\mathbb{T}} \theta (x) dx \sum_{k=0 }^\infty 2^{-2\alpha k} (\sin(2^{k } \pi y))^2,
\end{equation}
which is equal to $0$ when 
\begin{equation}
\int_{\mathbb{T}} \theta (x) dx =0.
\end{equation}
This proves point 1 of Proposition \ref{th1}. To prove point 2 one introduces the sequence $y_n=2^{-n}$ (which is converging to 0, as $n \rightarrow \infty$)  and consider the expression
\begin{equation}
\frac{1}{ y_n } U_{RR} (y_n;\theta)=2^{n-1}   \sum_{k =0}^{\infty} 2^{-2\alpha k} (\sin(2^{k } \pi 2^{-n}))^2.
\end{equation}
We first observe that $\sin(2^{k } \pi 2^{-n}) = 0$ for $k \geq n$. Therefore, the above sum is actually given by
\begin{equation}
\frac{1}{ y_n } U_{RR} (y_n;\theta)=2^{n-1}   \sum_{k =0}^{n-1} 2^{-2\alpha k} (\sin(2^{k -n} \pi ))^2.
\end{equation}
Now we observe that for $0 \leq k \leq n-1$ we have that $0 \leq 2^{k -n} \pi \leq \frac{\pi}{2}$. We recall that for $x \in \big[ 0, \frac{\pi}{2} \big]$ it holds that $\sin (x) \geq \frac{2}{\pi } x$. Applying this inequality to the series above gives
\begin{align*}
\frac{1}{ y_n } U_{RR} (y_n;\theta) &\geq 2^{n-1} \sum_{k =0}^{n-1} 2^{-2\alpha k} \bigg( \frac{2}{\pi} 2^{k -n} \pi \bigg)^2 = 2^{-n+1} \sum_{k =0}^{n-1} 2^{2k (1 - \alpha)} \\
&= 2^{-n+1} \frac{2^{2n(1-\alpha)} - 1}{2^{2(1-\alpha)} -1} = \frac{2}{2^{2(1-\alpha)} -1} \big( 2^{n(1-2\alpha)} - 2^{-n} \big).
\end{align*}
From the above we conclude that
\begin{equation}
\liminf_{n \rightarrow \infty} \frac{1}{ y_n } U_{RR} (y_n;\theta) \begin{cases}
= \infty \quad \text{if } 0 < \alpha < \frac{1}{2}, \\
\geq 2 \quad \text{if } \alpha = \frac{1}{2}, \\
\geq 0 \quad \text{if } \alpha > \frac{1}{2}.
\end{cases}
\end{equation}
Observing that 
\begin{equation*}
U(y,\theta) = U_{NR}(y,\theta) +  U_{R}(y,\theta) =  U_{NR}(y,\theta) + U_{RNR}(y,\theta) + U_{RR}(y,\theta)
\end{equation*}
completes the proof of point 2 of Proposition \ref{th1}. Since for $\alpha >\frac 12$ we have that $\lim 2^{n(1-2\alpha)}$ goes to $0$ with $n\rightarrow \infty$ point 2 is not in contradiction with point 1 of Proposition \ref{th1} and Lemma \ref{equivalencelemma}.
\end{proof}

\end{document}
